\chapter{Complexity — Meyer Tableau}
%%% generated by the blueprint pipeline from TCSlib.Complexity.CircuitComplexity.MeyerTableau
\section{Overview}

Declarations of the Lean module \texttt{TCSlib.Complexity.CircuitComplexity.MeyerTableau}.

\section{Declarations}

\begin{definition}[Shifted input symbol lookup]
\lean{Complexity.Meyer.xhat}
\label{Complexity.Meyer.xhat}
\leanok
Let $x$ be a finite Boolean string and $z\in\mathbb{Z}$. If $z\ge 1$, return the symbol
of $x$ at zero-based position $z-1$, if that position exists; otherwise, and whenever
$z<1$, return no symbol.
\end{definition}

\begin{definition}[Integer with a prescribed sign]
\lean{Complexity.Meyer.sgnZ}
\label{Complexity.Meyer.sgnZ}
\leanok
For a Boolean value $b$ and $a\in\mathbb{N}$, return $-a$ if $b$ is true and $a$ if $b$
is false.
\end{definition}

\begin{definition}[Work-tape symbol at a relative position]
\lean{Complexity.Meyer.RW}
\label{Complexity.Meyer.RW}
\leanok
\uses{Turing.Cfg}
Let $c$ be a configuration of a finite Turing machine, let $h$ be one of its work tapes,
and let $d\in\mathbb{Z}$. The value is the symbol, if any, stored on tape $h$ at
position $d$ relative to that tape’s head.
\end{definition}

\begin{definition}[Input symbol at a relative position]
\lean{Complexity.Meyer.RI}
\label{Complexity.Meyer.RI}
\leanok
\uses{Complexity.Meyer.xhat, Turing.Cfg}
Let $c$ be a configuration relative to a finite Boolean input $x$, and let
$d\in\mathbb{Z}$. The value is the input symbol, if any, at position $d$ relative to the
input head, using the input indexing in which the first input symbol has position $1$.
\end{definition}

\begin{theorem}[Input symbol at the current head position]
\lean{Complexity.Meyer.inputSymbol_eq_xhat}
\label{Complexity.Meyer.inputSymbol_eq_xhat}
\leanok
\uses{Complexity.Meyer.RI, Complexity.Meyer.xhat, Turing.Cfg, Turing.Cfg.inputSymbol, Turing.FinTM.inputSymbol_at}
\difficulty{4}
Let $M$ be a finite Turing machine over the Boolean alphabet, let $x$ be a finite
Boolean input, and let $c$ be a configuration of $M$ on $x$. The symbol read by the
input head of $c$ is the symbol of $x$ at the head's position, with no symbol returned
when that position is a boundary blank; equivalently, it is the input symbol at the
integer coordinate given by the input-head position.
\end{theorem}

\begin{theorem}[Current Work Tape Read]
\lean{Complexity.Meyer.workTapeSymbols_eq_RW}
\label{Complexity.Meyer.workTapeSymbols_eq_RW}
\leanok
\uses{Complexity.Meyer.RW, Turing.Cfg, Turing.Cfg.workTapeSymbols}
\difficulty{1}
For a configuration $c$ of a finite-tape Turing machine and a work-tape index $h$, the
symbol read by tape $h$ is the symbol stored at the head position; equivalently, it is
the symbol at offset $0$ from that head.
\end{theorem}

\begin{theorem}[Iterated rightward head movement]
\lean{Complexity.Meyer.iterate_move_pos}
\label{Complexity.Meyer.iterate_move_pos}
\leanok
\uses{Turing.moveInputPos, Turing.moveInputPos_pos_of_ne_right, Turing.moveInputPos_rightBoundary}
\difficulty{4}
Let $n,a\in\mathbb{N}$ and let $p$ be an input-head position among the $n+2$ positions.
After applying the positive one-step head movement $a$ times, the resulting position is
$\min(p+a,n+1)$.
\end{theorem}

\begin{theorem}[Repeated left moves of the input head]
\lean{Complexity.Meyer.iterate_move_neg}
\label{Complexity.Meyer.iterate_move_neg}
\leanok
\uses{Turing.moveInputPos, Turing.moveInputPos_leftBoundary, Turing.moveInputPos_neg_of_ne_left}
\difficulty{4}
Let $n\in\mathbb{N}$, let $p$ be one of the $n+2$ input-head positions, and let
$a\in\mathbb{N}$. After $a$ successive left moves, the head position is $p-a$, with
subtraction understood in the natural numbers and hence truncated at zero.
\end{theorem}

\begin{theorem}[Repeated input-head movement]
\lean{Complexity.Meyer.wstep_none_iterate}
\label{Complexity.Meyer.wstep_none_iterate}
\leanok
\uses{Complexity.Meyer.walkDir, Complexity.Meyer.wstep, Turing.Cfg, Turing.moveInputPos}
\difficulty{4}
Let $M$ be a finite-state Boolean Turing machine, let $x$ be a Boolean word, and let $c$
be a configuration of $M$ on input $x$. For any direction flag $neg$, blank flag $bl$,
and natural number $a$, applying the selector-directed step $a$ times yields the
configuration obtained from $c$ by applying the corresponding clamped input-head move
$a$ times to its input-head position; all other components of $c$ are unchanged.
\end{theorem}

\begin{theorem}[Input symbol after repeated walks]
\lean{Complexity.Meyer.walked_input}
\label{Complexity.Meyer.walked_input}
\leanok
\uses{Complexity.Meyer.RI, Complexity.Meyer.inputSymbol_eq_xhat, Complexity.Meyer.iterate_move_neg, Complexity.Meyer.iterate_move_pos, Complexity.Meyer.sgnZ, Complexity.Meyer.walkDir, Complexity.Meyer.wstep, Complexity.Meyer.wstep_none_iterate, Complexity.Meyer.xhat, Turing.Cfg, Turing.Cfg.inputSymbol}
\difficulty{5}
Let $M$ be a finite-state Turing machine with Boolean input symbols, let $x$ be a
Boolean word, and let $c$ be a configuration of $M$ on $x$. For any direction flag
$neg$, Boolean flag $bl$, and natural number $a$, after $a$ repetitions of the simulated
input-head walk in direction $neg$, the scanned input symbol equals the symbol of $x$ at
the original input-head position shifted by $-a$ when $neg$ is true and by $a$
otherwise, with positions outside $x$ read as blank.
\end{theorem}

\begin{theorem}[Repeated work-tape head movement]
\lean{Complexity.Meyer.wstep_some_iterate}
\label{Complexity.Meyer.wstep_some_iterate}
\leanok
\uses{Complexity.Meyer.sgnZ, Complexity.Meyer.walkDir, Complexity.Meyer.wstep, Turing.Cfg}
\difficulty{6}
Let $M$ be a finite-state Boolean Turing machine, let $x$ be a Boolean word, let $c$ be
a configuration of $M$ on input $x$, and choose a work tape $h$. For any direction flag
$neg$, blank flag $bl$, and natural number $a$, applying the selector-directed step $a$
times leaves every component of $c$ unchanged except the head position on tape $h$,
which becomes its original position plus $a$ when $neg$ is false and minus $a$ when
$neg$ is true.
\end{theorem}

\begin{theorem}[Work-tape symbol after repeated walks]
\lean{Complexity.Meyer.walked_work}
\label{Complexity.Meyer.walked_work}
\leanok
\uses{Complexity.Meyer.RW, Complexity.Meyer.sgnZ, Complexity.Meyer.wstep, Complexity.Meyer.wstep_some_iterate, Turing.Cfg, Turing.Cfg.workTapeSymbols}
\difficulty{3}
Let $M$ be a finite-state Turing machine with Boolean symbols, let $x$ be a Boolean
word, let $h$ index one of its work tapes, and let $c$ be a configuration of $M$ on $x$.
For any direction flag $neg$, Boolean flag $bl$, and natural number $a$, after $a$
repetitions of the simulated walk on tape $h$ in direction $neg$, the symbol scanned on
that tape is the original tape content at the original head position shifted by $-a$
when $neg$ is true and by $a$ otherwise.
\end{theorem}

\begin{theorem}[Work-tape answers after signed movement]
\lean{Complexity.Meyer.tabAns_sym_some}
\label{Complexity.Meyer.tabAns_sym_some}
\leanok
\uses{Complexity.Meyer.RW, Complexity.Meyer.selAnswer, Complexity.Meyer.sgnZ, Complexity.Meyer.symBit, Complexity.Meyer.tabAns, Complexity.Meyer.walked_work, Turing.Cfg}
\difficulty{1}
Let $M$ be a finite-tape machine over Boolean symbols, let $x$ be its input, and let $c$
be a configuration on $x$. For a work-tape index $h$, Boolean direction $\mathrm{neg}$,
Boolean blank flag $\mathrm{bl}$, and natural number $a$, the table answer for tape $h$
after $a$ simulated moves in that direction is the blank-symbol test applied to the
symbol at displacement $(-1)^{\mathrm{neg}}a$ from the current head position on tape
$h$.
\end{theorem}

\begin{theorem}[Input-tape answers after signed movement]
\lean{Complexity.Meyer.tabAns_sym_none}
\label{Complexity.Meyer.tabAns_sym_none}
\leanok
\uses{Complexity.Meyer.RI, Complexity.Meyer.selAnswer, Complexity.Meyer.sgnZ, Complexity.Meyer.symBit, Complexity.Meyer.tabAns, Complexity.Meyer.walked_input, Turing.Cfg}
\difficulty{1}
Let $M$ be a finite-tape machine over Boolean symbols, let $x$ be its input, and let $c$
be a configuration on $x$. For Boolean values $\mathrm{neg}$ and $\mathrm{bl}$ and a
natural number $a$, the table answer for the input-tape symbol selector, after $a$
simulated moves in direction $\mathrm{neg}$, is the blank-symbol test applied to the
input symbol at signed offset $(-1)^{\mathrm{neg}}a$ from the current input-head
position. The offset is $-a$ when $\mathrm{neg}$ is true and $a$ otherwise.
\end{theorem}

\begin{theorem}[State-query answers ignore the walk count]
\lean{Complexity.Meyer.tabAns_st}
\label{Complexity.Meyer.tabAns_st}
\leanok
\uses{Complexity.Meyer.selAnswer, Complexity.Meyer.tabAns, Complexity.TimeHierarchy.OutReg, Complexity.TimeHierarchy.OutReg.ofList, Turing.Cfg}
\difficulty{1}
Let $M$ be a finite-tape machine over Boolean symbols, and let $c$ be one of its
configurations. For every optional state $\sigma$, output-register value $r_0$, and
natural number $a$, the table answer for the state selector $(\sigma,r_0)$ is true
exactly when the configuration's state is $\sigma$ and its output, summarized as empty,
a single bit, or many bits, is $r_0$; this answer is independent of $a$.
\end{theorem}

\begin{definition}[Read symbols for a machine]
\lean{Complexity.Meyer.Rd}
\label{Complexity.Meyer.Rd}
\leanok
For a finite-tape Turing machine $M$ with Boolean input symbols, let $R_d(M)$ be the
product of an optional Boolean and a function assigning an optional Boolean to each of
the machine's work tapes.
\end{definition}

\begin{definition}[Symbols read in a configuration]
\lean{Complexity.Meyer.reads}
\label{Complexity.Meyer.reads}
\leanok
\uses{Complexity.Meyer.RI, Complexity.Meyer.RW, Complexity.Meyer.Rd, Turing.Cfg}
For a finite-tape Turing machine $M$ with Boolean input symbols and a configuration $c$
on input $x$, let the read data be the pair consisting of the symbol at input-tape
position $0$ and, for each work tape, the symbol at position $0$. Each symbol is
optional, with absence representing a blank cell.
\end{definition}

\begin{definition}[State and output summary]
\lean{Complexity.Meyer.SS}
\label{Complexity.Meyer.SS}
\leanok
\uses{Complexity.TimeHierarchy.OutReg, Complexity.TimeHierarchy.OutReg.ofList, Turing.Cfg}
For a finite-tape Turing machine $M$ with Boolean input symbols and a configuration $c$,
define its state-output summary to consist of its optional internal state and the output
register obtained from its output tape contents.
\end{definition}

\begin{definition}[State and output transition rule]
\lean{Complexity.Meyer.ruleS}
\label{Complexity.Meyer.ruleS}
\leanok
\uses{Complexity.Meyer.Rd, Complexity.TimeHierarchy.OutReg, Complexity.TimeHierarchy.OutReg.push}
For a finite-tape Turing machine $M$ with Boolean input symbols, a state-output summary
$(q,r)$, and read data $\rho$, define the updated summary as follows. If $q$ is absent,
retain $(q,r)$; if $q$ is present, apply the machine transition using the input and
work-tape symbols in $\rho$, take its resulting state, and push its output onto $r$.
\end{definition}

\begin{definition}[Work-tape symbol update]
\lean{Complexity.Meyer.wsym}
\label{Complexity.Meyer.wsym}
\leanok
Given an optional optional Boolean $w$ and an optional Boolean $r$, define the updated
symbol to be $r$ when $w$ is absent, and to be the inner optional Boolean when $w$ is
present.
\end{definition}

\begin{definition}[Work-tape cell update rule]
\lean{Complexity.Meyer.ruleW}
\label{Complexity.Meyer.ruleW}
\leanok
\uses{Complexity.Meyer.Rd, Complexity.Meyer.wsym}
For a finite-tape Turing machine $M$ with Boolean input symbols, an optional state $q$,
read data $\rho$, a work tape $h$, an integer cell offset $d$, and tape contents
$f:\mathbb{Z}\to\mathrm{Option}(\mathrm{Bool})$, define the updated cell symbol as
follows. If $q$ is absent, use $f(d)$. If $q$ is present, let the transition's movement
on tape $h$ be $m$ and its write instruction be $w$: when $d+m=0$, use the write
instruction applied to the symbol read on tape $h$; otherwise use $f(d+m)$.
\end{definition}

\begin{definition}[Input-tape cell update rule]
\lean{Complexity.Meyer.ruleI}
\label{Complexity.Meyer.ruleI}
\leanok
\uses{Complexity.Meyer.Rd}
For a finite-tape Turing machine $M$ with Boolean input symbols, an optional state $q$,
read data $\rho$, an integer offset $d$, and input-tape contents
$f:\mathbb{Z}\to\mathrm{Option}(\mathrm{Bool})$, define the updated cell symbol as
follows. If $q$ is absent, use $f(d)$. If $q$ is present, retain $f(d)$ when the input
symbol read is blank and the transition's input-tape movement leads to a blank cell;
otherwise use $f(d+m)$, where $m$ is that transition's input-tape movement.
\end{definition}

\begin{theorem}[Transition Step from a Nonhalting Configuration]
\lean{Complexity.Meyer.step_eq}
\label{Complexity.Meyer.step_eq}
\leanok
\uses{Complexity.Meyer.inputSymbol_eq_xhat, Complexity.Meyer.reads, Complexity.Meyer.workTapeSymbols_eq_RW, Turing.Action.apply, Turing.Cfg, Turing.Cfg.workTapeSymbols, Turing.MultiTapeTM.step}
\difficulty{4}
Let $M$ be a finite Turing machine, let $x$ be an input string, and let $c$ be a
configuration of $M$ on $x$. If the state of $c$ is $q$, then one step from $c$ is
obtained by applying the transition action prescribed by $M$ at state $q$, using the
symbols currently read on the input and work tapes.
\end{theorem}

\begin{theorem}[State and output summary after one step]
\lean{Complexity.Meyer.SS_step}
\label{Complexity.Meyer.SS_step}
\leanok
\uses{Complexity.Meyer.SS, Complexity.Meyer.reads, Complexity.Meyer.ruleS, Complexity.Meyer.step_eq, Complexity.TimeHierarchy.OutReg.ofList_push, Turing.Action.apply, Turing.Cfg, Turing.MultiTapeTM.step}
\difficulty{4}
Let $M$ be a finite-state Turing machine with $k$ work tapes, let $x$ be a Boolean input
word, and let $c$ be a configuration of $M$ on $x$. The state and output summary after
one step is obtained by applying the summary update rule to the current state and output
summary of $c$ together with the symbols read by the input and work-tape heads in $c$.
\end{theorem}

\begin{theorem}[Work-tape contents after one step]
\lean{Complexity.Meyer.RW_step}
\label{Complexity.Meyer.RW_step}
\leanok
\uses{Complexity.Meyer.RW, Complexity.Meyer.reads, Complexity.Meyer.ruleW, Complexity.Meyer.step_eq, Complexity.Meyer.wsym, Turing.Action.apply, Turing.Cfg, Turing.MultiTapeTM.step, Turing.MultiTapeTM.step_of_halt}
\difficulty{5}
Let $M$ be a finite-state Turing machine with $k$ work tapes, let $x$ be a Boolean input
word, and let $c$ be a configuration of $M$ on $x$. For each work tape $h$ and integer
offset $d$ from its head, the symbol at that offset after one machine step is exactly
the value prescribed by the work-tape update rule for the current state, the symbols
read in $c$, and the offset $d$, with the prior tape contents supplied as the rule's
default.
\end{theorem}

\begin{theorem}[When the shifted word is blank]
\lean{Complexity.Meyer.xhat_eq_none}
\label{Complexity.Meyer.xhat_eq_none}
\leanok
\uses{Complexity.Meyer.xhat}
\difficulty{3}
Let $x$ be a Boolean word and $z$ an integer. The shifted word has no symbol at position
$z$ exactly when $z<1$ or $|x|<z$.
\end{theorem}

\begin{theorem}[Input Read After One Step]
\lean{Complexity.Meyer.RI_step}
\label{Complexity.Meyer.RI_step}
\leanok
\uses{Complexity.Meyer.RI, Complexity.Meyer.reads, Complexity.Meyer.ruleI, Complexity.Meyer.step_eq, Complexity.Meyer.xhat, Complexity.Meyer.xhat_eq_none, Turing.Action.apply, Turing.Cfg, Turing.MultiTapeTM.step, Turing.MultiTapeTM.step_of_halt, Turing.moveInputPos, Turing.moveInputPos_leftBoundary, Turing.moveInputPos_neg_of_ne_left, Turing.moveInputPos_pos_of_ne_right, Turing.moveInputPos_rightBoundary}
\difficulty{6}
Let $M$ be a finite-tape Turing machine with Boolean symbols, let $x$ be a finite
Boolean input word, let $c$ be a configuration of $M$ on $x$, and let $d\in\mathbb{Z}$.
Write $f(e)$ for the input symbol seen from $c$ at offset $e$, and let $\rho$ record the
input and work-tape symbols read in $c$. If $c$ is halted, the input read after one step
at offset $d$ is $f(d)$. If $c$ is active, let $m$ be the input-head movement prescribed
by $M$ from its state and reads $\rho$. The read after one step is $f(d)$ when the input
read in $\rho$ is absent and $f(m)$ is also absent; in every other case it is $f(d+m)$.
\end{theorem}

\begin{theorem}[Initial state and output summary]
\lean{Complexity.Meyer.SS_init}
\label{Complexity.Meyer.SS_init}
\leanok
\uses{Complexity.Meyer.SS, Turing.MultiTapeTM.initCfg}
\difficulty{1}
Let $M$ be a finite-state Turing machine with initial state $q_0$, and let $x$ be a
Boolean input word. The state and output summary of its initial configuration on $x$ is
$(q_0,	ext{empty})$.
\end{theorem}

\begin{theorem}[Initial work-tape contents]
\lean{Complexity.Meyer.RW_init}
\label{Complexity.Meyer.RW_init}
\leanok
\uses{Complexity.Meyer.RW, Turing.MultiTapeTM.initCfg}
\difficulty{1}
Let $M$ be a finite-state Turing machine with $k$ work tapes and let $x$ be a Boolean
input word. In the initial configuration of $M$ on $x$, for every work tape $h$ and
every integer offset $d$ from its head, the scanned cell contains no symbol.
\end{theorem}

\begin{theorem}[Input Position in the Initial Configuration]
\lean{Complexity.Meyer.RI_init}
\label{Complexity.Meyer.RI_init}
\leanok
\uses{Complexity.Meyer.RI, Complexity.Meyer.xhat, Turing.MultiTapeTM.initCfg}
\difficulty{1}
Let $M$ be a finite-tape Turing machine with Boolean symbols and let $x$ be a finite
Boolean input word. Define the extended input lookup $\widehat{x}(z)$ for
$z\in\mathbb{Z}$ to be the $(z-1)$-st entry of $x$ when $z\geq1$ and that index is in
range, and absent otherwise. In the initial configuration of $M$ on $x$, the input
symbol seen at offset $d\in\mathbb{Z}$ is $\widehat{x}(1+d)$.
\end{theorem}

\begin{theorem}[Length of a bounded binary word]
\lean{Complexity.Meyer.length_enumWord}
\label{Complexity.Meyer.length_enumWord}
\leanok
\uses{Complexity.enumWord}
\difficulty{4}
For all natural numbers $w$ and $v$, the width-$w$ little-endian binary word of $v$ has
length $w$.
\end{theorem}

\begin{theorem}[Value of a Fixed-Width Binary Word]
\lean{Complexity.Meyer.ctrVal_enumWord}
\label{Complexity.Meyer.ctrVal_enumWord}
\leanok
\uses{Complexity.TimeHierarchy.ctrVal, Complexity.enumWord}
\difficulty{5}
For natural numbers $w,v$ with $v<2^w$, the binary value of the width-$w$ little-endian
binary word of $v$ is $v$.
\end{theorem}

\begin{theorem}[Binary word recovers its integer value]
\lean{Complexity.Meyer.enumWord_ctrVal}
\label{Complexity.Meyer.enumWord_ctrVal}
\leanok
\uses{Complexity.TimeHierarchy.ctrVal, Complexity.enumWord}
\difficulty{5}
For every finite Boolean word $u$, the width-$|u|$ little-endian binary word of the
integer represented by $u$ is exactly $u$, where the represented integer is obtained by
summing the bits with successive powers of $2$ as weights.
\end{theorem}

\begin{theorem}[Zero binary word]
\lean{Complexity.Meyer.enumWord_zero}
\label{Complexity.Meyer.enumWord_zero}
\leanok
\uses{Complexity.enumWord}
\difficulty{4}
For every natural number $w$, the width-$w$ little-endian binary word of $0$ is the word
of length $w$ consisting entirely of $0$s.
\end{theorem}

\begin{theorem}[All-ones binary word]
\lean{Complexity.Meyer.enumWord_ones}
\label{Complexity.Meyer.enumWord_ones}
\leanok
\uses{Complexity.enumWord}
\difficulty{4}
For every natural number $w$, the width-$w$ little-endian binary word of $2^w-1$ is the
word of length $w$ consisting entirely of $1$s.
\end{theorem}

\begin{theorem}[Incrementing an enumerated binary word]
\lean{Complexity.Meyer.incFixed_enumWord}
\label{Complexity.Meyer.incFixed_enumWord}
\leanok
\uses{Complexity.enumWord, Turing.incFixed}
\difficulty{5}
For natural numbers $w$ and $v$ satisfying $v+1<2^w$, applying fixed-width binary
increment to the width-$w$ little-endian binary word of $v$ yields the width-$w$
little-endian binary word of $v+1$.
\end{theorem}

\begin{theorem}[Overflow of fixed-width increment]
\lean{Complexity.Meyer.incFixed_eq_none_iff}
\label{Complexity.Meyer.incFixed_eq_none_iff}
\leanok
\uses{Turing.incFixed}
\difficulty{4}
For every finite Boolean word $u$, the fixed-width increment overflows exactly when
every bit of $u$ is true; equivalently, its result is none if and only if $u$ is the
word of length $|u|$ consisting entirely of true bits.
\end{theorem}

\begin{theorem}[Appending Zeros Preserves Binary Value]
\lean{Complexity.Meyer.ctrVal_append_zeros}
\label{Complexity.Meyer.ctrVal_append_zeros}
\leanok
\uses{Complexity.TimeHierarchy.ctrVal}
\difficulty{4}
For every finite Boolean word $l$ and every $k\in\mathbb{N}$, appending $k$ false bits
to $l$ leaves its binary value unchanged.
\end{theorem}

\begin{theorem}[Fixed-width binary representation]
\lean{Complexity.Meyer.bits_pad}
\label{Complexity.Meyer.bits_pad}
\leanok
\uses{Complexity.Meyer.ctrVal_append_zeros, Complexity.Meyer.enumWord_ctrVal, Complexity.TimeHierarchy.ctrVal_bits, Complexity.enumWord}
\difficulty{4}
For natural numbers $w,j$ with $j<2^w$, the binary digits of $j$, listed least
significant first and followed by enough zeroes to give a word of length $w$, equal the
width-$w$ little-endian binary word of $j$.
\end{theorem}

\begin{theorem}[Query Evaluation]
\lean{Complexity.Meyer.tabFun_query}
\label{Complexity.Meyer.tabFun_query}
\leanok
\uses{Complexity.Meyer.Sel, Complexity.Meyer.encodeSel, Complexity.Meyer.parseQ_query, Complexity.Meyer.query, Complexity.Meyer.tabAns, Complexity.Meyer.tabFun, Complexity.TimeHierarchy.ctrVal, Turing.MultiTapeTM.initCfg, Turing.MultiTapeTM.runFrom}
\difficulty{1}
For a Boolean Turing machine $M$, a selector $s$, and Boolean lists $t$, $a$, and $x$,
the tabulation function on the query encoding of $s,t,a,x$ equals the selector’s answer
on the configuration obtained by running $M$ from its initial configuration on input $x$
for the integer represented by $t$, with lookahead $a$.
\end{theorem}
